% Blank variant of src/architecture/decisions/templates/adr-cloud-platform-template.tex. Unapproved working document.
% --------------------------------------------------------------------------
% File: src/architecture/decisions/templates/adr-cloud-platform-template.tex
% Style: style v1.0 (tooling/latex/style.sty)
% Engine: pdfLaTeX through tooling/scripts/latex_build.py
% Purpose: ready-to-fill blank record; not an approved decision
% --------------------------------------------------------------------------
\documentclass[11pt,letterpaper]{article}
\usepackage{technical-design-spec}
\renewcommand{\DocTitle}{ADR: Cloud Platform and Landing Zone}
\renewcommand{\DocSubtitle}{Fillable working record and evidence trail}
\renewcommand{\DocVersion}{v1.0}
\renewcommand{\DocStatus}{Blank}
\renewcommand{\DocOwner}{Architecture Governance}
\renewcommand{\DocAuthor}{Architecture Documentation Library}
\renewcommand{\DocDate}{2026-10-28}
\renewcommand{\DocSummary}{An unfilled, domain-specific architecture decision record for cloud platform and landing zone. Inherits the central decision structure and adds specialized design evidence.}
\begin{document}
\makedoctocpage


\section{Record identity and decision authority}
\noindent\textbf{Status:} \ADRStatus{Proposed}\quad
\textbf{ADR ID:} \emph{[ADR-0000]}\quad
\textbf{Date:} \emph{[YYYY-MM-DD]}

\begin{description}
\item[Decision title:] \emph{[To complete]}
\item[Entity of interest:] \emph{[To complete]}
\item[Scope and baseline:] \emph{[To complete]}
\item[Owner and accountable approver:] \emph{[To complete]}
\item[Consulted stakeholders:] \emph{[To complete]}
\item[Related records:] \emph{[To complete]}
\item[Review and expiry trigger:] \emph{[To complete]}
\end{description}

\begin{notebox}[Unapproved working record]
Complete every required field. Acceptance requires recorded decision
authority and evidence; this Blank is never an approved decision.
\end{notebox}

\section{Context and architectural significance}
\begin{description}
\item[Problem and forces:] \emph{[To complete]}
\item[Stakeholders and concerns:] \emph{[To complete]}
\item[Drivers and constraints:] \emph{[To complete]}
\item[Boundary and assumptions:] \emph{[To complete]}
\item[Decision deadline:] \emph{[To complete]}
\end{description}

\section{Options considered and evidence}
Include the status quo as an explicit alternative where practical. Evaluate
at least two feasible options, or explain why only one remains viable.
\begin{tabularx}{\linewidth}{@{}p{0.18\linewidth}X X@{}}
\toprule
\textbf{Option} & \textbf{Benefits and supporting evidence} & \textbf{Costs, risks and reason for disposition}\\
\midrule
A: \emph{[To complete]} & \emph{[To complete]} & \emph{[To complete]}\\
B: \emph{[To complete]} & \emph{[To complete]} & \emph{[To complete]}\\
C: \emph{[To complete]} & \emph{[To complete]} & \emph{[To complete]}\\
\bottomrule
\end{tabularx}

\noindent\textbf{Evaluation method:} \emph{[To complete]}

\section{Decision, rationale and non-goals}
\begin{description}
\item[We choose:] \emph{[To complete]}
\item[Because:] \emph{[To complete]}
\item[Non-goals:] \emph{[To complete]}
\item[Invariants and exceptions:] \emph{[To complete]}
\item[Dependencies:] \emph{[To complete]}
\end{description}

\section{Domain-specific decision inputs}
Complete the following prompts with referenceable designs, measurements,
threat assessments or operational constraints; mark inapplicable prompts
explicitly rather than removing necessary decision context.
\begin{itemize}
\item Placement model: provider(s), region(s), residency, tenancy, environment isolation and managed-service responsibility boundaries.
\item Account/subscription hierarchy: organizational units, guardrails, policy inheritance, delegated administration and separation of duties.
\item Network architecture: ingress/egress, private connectivity, DNS, segmentation, endpoints, hybrid integration and threat boundaries.
\item Identity and key management: workforce/workload federation, least-privilege roles, key ownership, rotation and break-glass access.
\item Platform lifecycle: baseline configuration, golden paths, patching, drift detection, shared services and decommissioning.
\end{itemize}
\section{Domain-specific alternatives and trade-offs}
Explain how each feasible option addresses the following criteria.
\begin{itemize}
\item Isolation versus operational overhead; single-region versus multi-region; self-managed versus managed service.
\item Provider concentration and exit cost; data-transfer cost, tagging, budgets, service quotas and demand forecasts.
\item Boundary ownership between central platform, security, application teams and providers.
\end{itemize}
\noindent\textbf{Domain comparison:} \emph{[To complete]}
\section{Domain-specific validation and review gates}
The proposed decision is ready for acceptance only after its applicable
tests have objective pass/fail thresholds, an owner and dated evidence.
\begin{itemize}
\item Infrastructure plan/policy-as-code checks and account-level least-privilege tests.
\item Network reachability tests (allowed and denied), encryption/key policy evidence, audit-log retention.
\item Cost model under expected and stress traffic, failure-domain test and Terraform state isolation verification.
\end{itemize}
\noindent\textbf{Related repository material:} Cloud mechanisms catalog; cloud security governance templates; workload-specific architecture descriptions.

\section{Consequences, risks and safeguards}
\begin{description}
\item[Positive consequences:] \emph{[To complete]}
\item[Negative consequences:] \emph{[To complete]}
\item[Security and privacy effects:] \emph{[To complete]}
\item[Operational and financial effects:] \emph{[To complete]}
\item[Unknowns and risk treatment:] \emph{[To complete]}
\end{description}

\section{Implementation, verification and reversibility}
\begin{tabularx}{\linewidth}{@{}p{0.22\linewidth}X@{}}
\toprule
\textbf{Control} & \textbf{Complete with concrete, reviewable evidence}\\
\midrule
Implementation owner and plan & \emph{[To complete]}\\
Acceptance tests & \emph{[To complete]}\\
Security / reliability gates & \emph{[To complete]}\\
Evidence and provenance & \emph{[To complete]}\\
Rollback or exit & \emph{[To complete]}\\
Post-deployment review & \emph{[To complete]}\\
\bottomrule
\end{tabularx}

\section{Architecture-description traceability}
A decision record contributes rationale to an architecture description; it
does not, by itself, constitute a conforming ISO/IEC/IEEE 42010:2022
architecture description.
\begin{description}
\item[Description and configuration baseline:] \emph{[To complete]}
\item[Stakeholders and concerns:] \emph{[To complete]}
\item[Viewpoints, views and model kinds:] \emph{[To complete]}
\item[Correspondence and consistency:] \emph{[To complete]}
\item[Requirements and evidence:] \emph{[To complete]}
\end{description}

\section{Review history and references}
\begin{tabularx}{\linewidth}{@{}p{0.16\linewidth}p{0.19\linewidth}p{0.19\linewidth}X@{}}
\toprule
\textbf{Date} & \textbf{Reviewer / role} & \textbf{Disposition} & \textbf{Evidence and requested changes}\\
\midrule
\emph{[date]} & \emph{[To complete]} & \emph{[To complete]} & \emph{[To complete]}\\
\bottomrule
\end{tabularx}

\noindent\textbf{Sources:} \emph{[To complete]}

\end{document}
